% =====================================================================================
% SUITES DU POINT AVEC CAROLINE HILLAIRET -- etat point par point au 13 aout 2026.
%
% OBJET : dire ce qui est integre, ce qui ne l'est pas, et POURQUOI dans les deux cas. Le
% document est concu pour etre lu en dix minutes avant un point, donc une ligne par demande.
%
% LE VOCABULAIRE DE STATUT PORTE L'ESSENTIEL DE L'INFORMATION, et il distingue cinq cas qu'on
% confond d'ordinaire sous « fait / pas fait » :
%   INTEGRE                 : fait, sous la forme demandee ;
%   INTEGRE AUTREMENT       : fait, mais la mesure a conduit ailleurs que l'intuition ;
%   MESURE PUIS ECARTE      : le calcul a ete fait ET il dit non. C'est le statut le plus
%                             defendable, et le plus facile a confondre avec « pas fait » ;
%   CHIFFRE, DECISION OUVERTE : le calcul est la, l'arbitrage revient a Kelian ;
%   NON TRAITE              : avec la raison, jamais sans.
% =====================================================================================
\documentclass[11pt,a4paper]{article}

\usepackage[utf8]{inputenc}
\usepackage[T1]{fontenc}
\usepackage{lmodern}
\usepackage[french]{babel}
\usepackage[margin=1.7cm]{geometry}
\usepackage{microtype}
\usepackage{array}
\usepackage{booktabs}
\usepackage{longtable}
\usepackage[table]{xcolor}
\usepackage[most]{tcolorbox}
\usepackage{fancyhdr}
\usepackage{lastpage}
% \euro vient du preambule du memoire ; ce document est autonome et doit le charger lui-meme.
\usepackage{eurosym}

\definecolor{navy}{HTML}{1F3864}
\definecolor{lightblue}{HTML}{D6E0F0}
\definecolor{vert}{HTML}{D9EAD3}
\definecolor{ambre}{HTML}{FCE5CD}
\definecolor{gris}{HTML}{EFEFEF}
\definecolor{rouge}{HTML}{F4CCCC}

\newtcolorbox{cle}[1][]{colback=lightblue!40,colframe=navy,boxrule=0.6pt,
  arc=2pt,left=8pt,right=8pt,top=6pt,bottom=6pt,#1}

\newcommand{\ok}{\cellcolor{vert}\footnotesize\textbf{INTÉGRÉ}}
\newcommand{\autrement}{\cellcolor{vert}\footnotesize\textbf{INTÉGRÉ}\newline\footnotesize\textbf{AUTREMENT}}
\newcommand{\ecarte}{\cellcolor{ambre}\footnotesize\textbf{MESURÉ PUIS}\newline\footnotesize\textbf{ÉCARTÉ}}
\newcommand{\ouvert}{\cellcolor{ambre}\footnotesize\textbf{CHIFFRÉ,}\newline\footnotesize\textbf{À TRANCHER}}
\newcommand{\non}{\cellcolor{rouge}\footnotesize\textbf{NON TRAITÉ}}
\newcommand{\toi}{\cellcolor{gris}\footnotesize\textbf{HORS}\newline\footnotesize\textbf{PÉRIMÈTRE}}

\pagestyle{fancy}
\fancyhf{}
\lfoot{\footnotesize Suites du point avec Caroline Hillairet --- 13 août 2026}
\rfoot{\footnotesize page \thepage/\pageref{LastPage}}
\renewcommand{\headrulewidth}{0pt}

\title{\vspace{-1.3cm}\color{navy}\bfseries Suites du point avec Caroline Hillairet\\[2pt]
  \large Ce qui est intégré, ce qui ne l'est pas, et pourquoi}
\author{Kélian Kaddouri}
\date{13 août 2026}

\begin{document}
\maketitle
\thispagestyle{fancy}
\vspace{-0.7cm}

\begin{cle}
\textbf{Comment lire ce tableau.} Deux statuts méritent d'être distingués de \og pas fait \fg,
et c'est le point principal du document. \textbf{Mesuré puis écarté} signifie que le calcul
demandé a bien été conduit et qu'il conclut contre la proposition : c'est un résultat, pas un
renoncement, et c'est le cas de la renormalisation par la taille. \textbf{Intégré autrement}
signifie que la demande est traitée mais que la mesure a conduit ailleurs que l'intuition qui la
motivait : c'est le cas de l'étage de modèle et de la sensibilité à la propagation. Les deux
seules lignes réellement ouvertes sont signalées comme telles.
\end{cle}

\vspace{0.3cm}
{\footnotesize
\renewcommand{\arraystretch}{1.25}
\begin{longtable}{@{}p{3.4cm} >{\raggedright\arraybackslash\hyphenpenalty=10000\exhyphenpenalty=10000}p{2.4cm} p{8.7cm} p{2.0cm}@{}}
\toprule
\textbf{Demande} & \textbf{Statut} & \textbf{Ce qui a été fait, et ce que cela donne} &
\textbf{Où le voir} \\
\midrule
\endfirsthead
\toprule
\textbf{Demande \emph{(suite)}} & \textbf{Statut} & \textbf{Ce qui a été fait} &
\textbf{Où} \\
\midrule
\endhead

Dynamique de $Y_{j,t}$, brownien pour le fond et sauts pour les chocs &
\ok &
Note de six pages écrite. Trois horloges distinguées, et les deux échelles sont
\emph{identifiées} par la surdispersion mesurée : $1{,}19$ sur les cellules firme-année,
$2{,}24$ sur la série annuelle détendancée, contre $1$ pour un brownien pur. La forme retenue est
un renouvellement, avec l'hypothèse déclarée plutôt que masquée. &
\texttt{note\_dyna\-mique\_Y} \\

Clarifier l'emploi du mot \og gravité \fg &
\autrement &
Balayage fait, et le résultat est inattendu : \textbf{le mémoire n'emploie nulle part
\og gravité \fg{} pour le caractère systémique}. La confusion venait des decks de juillet, qui
sont archivés. \textbf{En revanche il y a une vraie collision dans le code}, entre \texttt{G\_BASE}
(le gain de propagation) et \texttt{GBASE} (l'échelon de sévérité) : deux objets sans rapport dont
les noms ne diffèrent que par un tiret bas. Le renommage touche 27 fichiers et le pipeline est
gelé : la distinction est verrouillée dans la table des notations en attendant l'arbitrage. &
ch.~notations \\

Test de martingalité &
\ok &
Vérifié : \textbf{le mémoire n'en contient aucun}, et il n'y en a jamais eu. Le mot a deux
emplois, tous deux légitimes et sous mesure physique : la forme de Dirichlet d'une chaîne de
Markov, et un brownien sans dérive pour un temps de premier passage. Deux paragraphes de
précaution le disent désormais explicitement, et le lien rhétorique qui suggérait un cadre
commun est corrigé : ce que les deux usages partagent est la mesure, pas la mathématique. &
annexe~D \\

Sensibilité des paramètres &
\autrement &
Section écrite, et le résultat va \emph{contre} l'inquiétude qui la motivait. \textbf{La
propagation est le levier le moins sensible des quatre} : $g$ de $0{,}5$ à $1{,}0$ déplace le
surcoût de $5\,517$ à $6\,579$~M\euro, quand le seuil de la loi de queue le déplace de $4\,862$ à
$12\,944$ et la fréquence de $4\,056$ à $8\,684$. La fragilité est dans l'ajustement de valeurs
extrêmes, pas dans la contagion. &
ch.~résultats \\

Sensibilité à la direction de la contagion &
\ok &
Table publiée pour la première fois : l'ignorance sur la direction coûte $336$~M\euro{} de
largeur au quart, $695$ à la moitié, $1\,083$ aux trois quarts et $1\,839$ à l'ignorance totale.
\textbf{L'ignorance se paye linéairement et son coût maximal est connu d'avance}, ce qui est plus
favorable qu'un intervalle de confiance ordinaire. &
ch.~résultats \\

Reporter le SCR avec un intervalle à $95\,\%$ &
\ouvert &
Chiffré sans rejouer aucun bootstrap, par la même approximation normale que le moteur. L'indice
de queue passe à $[0{,}2487\,;0{,}8765]$, le capital à $[2\,727\,;34\,942]$~M\euro. Deux constats :
le passage de 90 à $95\,\%$ ne déplace \textbf{aucune calibration}, seulement les bornes
publiées ; et il ne corrige pas la sous-couverture, qui reste un déficit constant de
$-3{,}2$~points ($86{,}8\,\%$ réels pour 90 nominaux, $91{,}8$ pour 95). \textbf{La bascule est un
choix de présentation, à trancher.} &
annexe~D \\

CTE / CVaR à côté de la VaR &
\ok &
Calculée au niveau agrégé, en \emph{diagnostic} et jamais comme mesure de couverture.
Résultat qui ferme la discussion : \textbf{la CTE à $95\,\%$ est moins prudente que la VaR à
$99{,}5\,\%$}, l'équivalence tombant à $\beta = 97{,}73\,\%$. Et l'argument décisif n'est pas une
proportion mais une existence : sous la calibration alternative $\hat\xi > 1$, l'espérance est
infinie et \textbf{la mesure n'existe pas}. Une couverture qui cesse d'être définie selon la
source ne peut pas porter un capital. &
annexe~D \\

Renormaliser proportionnellement au SCR de marché &
\ecarte &
\textbf{Testé, et la mesure conclut contre.} La proportionnalité affirme une élasticité de $1$ à
la taille ; l'élasticité mesurée vaut $0{,}204$. Une petite entité se voit attribuer $0{,}633$
contre $0{,}100$ sous proportionnalité, une grande $1{,}621$ contre $10{,}00$. \textbf{Les deux
approches se trompent en sens opposés} : la proportionnalité sous-estime une petite entité, le
modèle la majore. Substituer l'une à l'autre changerait le signe du défaut sans le déclarer, donc
la proportionnalité reste le \emph{repère} contre lequel la méthode se distingue, pas une
méthode. &
ch.~résultats \\

Chiffrer les trois étages, puis retirer le troisième &
\autrement &
Chiffré, et \textbf{le résultat va contre l'intuition de l'arbitrage} : l'étage de modèle est le
plus gros des trois, facteur $9{,}21$ sur l'axe de la famille de queue contre $8{,}93$ pour le
paramètre et $1{,}27$ pour l'identification. Le retirer rétrécirait donc l'affichage plus que
tout autre retrait, alors que l'incertitude ne bougerait pas. Compromis retenu : il \textbf{sort
de la bande reportée} mais devient un \emph{axe de sensibilité déclaré}, et reste publié deux
fois. Motif écrit : les deux premiers étages sont statistiques donc portables par un intervalle,
le troisième est épistémique et un choix de famille ne se moyenne pas. &
ch.~robustesse \\

Méthode de Cooke en annexe, questionnaire peaufiné &
\ok &
Annexe écrite. Elle ne présente pas un protocole à venir mais \textbf{un protocole préparé et non
exécuté}, avec le chiffrage de la raison, ce qui est un actif devant un jury qui demanderait
pourquoi ce modèle ne recourt pas au jugement d'expert. Le questionnaire a été repris le 13 août :
une graine \textbf{fausse} en a été retirée (elle annonçait $\xi = 0{,}90$, valeur posée du pire
cas, quand la calibration donne $0{,}5954$ ; dans la méthode de Cooke une graine fausse
\emph{inverse} les poids), et la convention de direction, qui était inversée, est alignée. &
annexe~E \\

Conformité aux piliers 1 et 2 protégeant partiellement le 3 &
\ok &
Mesuré, et la réponse tient en trois temps. \textbf{Dire que la dépendance manque serait faux} :
elle existe déjà, de corrélation $0{,}462$ entre deux piliers quelconques. Ce qui manque est sa
\emph{structure}, échangeable au lieu d'être propre à certains triplets. \textbf{L'omission ne
coûte rien} : en ajoutant un surcroît de corrélation aux seules paires concernées, la loi des
configurations bouge nettement mais le capital espéré ne bouge que de $2$~M\euro, soit
$0{,}02\,\%$. \textbf{Et la raison est structurelle}, donc l'invariance vaut pour toute structure
à marges fixées : le capital s'ajuste par une forme additive en indicateurs de pilier à
$R^2 = 0{,}9945$, et l'espérance d'une fonction additive ne dépend que des marges. &
annexe~D \\

Défaillances simultanées, additivité des coûts &
\ok &
La prémisse est à corriger d'abord : \textbf{ce n'est pas un cas laissé de côté, c'est la sortie
du modèle}, et à l'état non conforme la défaillance multiple est \emph{majoritaire},
$62{,}31\,\%$ des sinistres contre $31{,}15\,\%$ à l'état conforme. Ensuite, \og additivité \fg{}
désigne \textbf{trois choses à trois étages} : hypothèse non testée sur les coûts d'un sinistre,
quasi-additivité \emph{mesurée} sur les piliers ($R^2 = 0{,}9945$), super-additivité
\emph{mesurée} sur les canaux ($+35\,\%$). Répondre \og le modèle est super-additif \fg{} serait
donc faux. Enfin l'hypothèse est bornée : un exposant de $0{,}70$ à $1{,}30$ déplace l'écart de
$9\,706$ à $20\,487$~M\euro. &
ch.~robustesse \\

Probabilités conditionnelles de séquence &
\autrement &
La voie a été suivie jusqu'au bout et se referme, pour trois raisons distinctes. Les séquences
étaient \textbf{déjà dans le corpus}, implicites : composer les arêtes donne $9$ chemins et
$4$ séquences distinctes. \textbf{La comparaison naturelle a un support vide} : P2 n'émet jamais,
$0$ fois sur $22$ transitions. \textbf{Le test n'a aucune puissance}, et pas faute d'échantillon :
il exige d'un même pilier qu'il soit atteint \emph{et} qu'il ait deux successeurs, et les deux
manques sont exclusifs ici. Acquis positif : la chaîne \textbf{n'est pas sans mémoire},
l'auto-évitement gonflant la loi du successeur de $1{,}211$ en moyenne. &
ch.~identifiabilité \\

Processus des dates d'arrivée, effet du regroupement &
\non &
\textbf{Non traité, et la raison est dans la donnée.} Le mécanisme manquant serait une saturation
de la capacité de réponse. Or le tableau des proxys fait apparaître que \textbf{le coût d'un
sinistre daté n'est jamais observé} dans les sources retenues : la chronologie de brèches porte
les dates, la conversion en montants porte les coûts, et aucune source ne porte les deux. Un
mécanisme de saturation ne serait donc calibrable par rien. La limite est déclarée plutôt
qu'estimée. &
annexe~D \\

Contacter Samuel Cywie &
\toi &
Les trois formulaires sont prêts et à jour au 13 août. L'envoi et le contact ne relèvent pas de
la préparation technique. &
--- \\

Lancer l'élicitation &
\ecarte &
\textbf{Abandonnée, sur décision assumée}, et le protocole en donne la raison chiffrée plutôt que
le calendrier : un panel faiblement calibré déplacerait le capital de $777$~M\euro{} \emph{sans
qu'aucun signal ne l'indique}, là où l'approche par bornes reste invariante. Une élicitation
faible introduirait donc une incertitude \textbf{non déclarable} en échange d'une ignorance
mesurée. Le coût de l'abandon est nommé : P1 et P4 restent à égalité, pour un regret maximal de
$1{,}5\,\%$. La décision est \textbf{réversible sans coût de développement}. &
annexe~E \\
\bottomrule
\end{longtable}
}

\vspace{0.15cm}
% RATTACHEMENT AU HARNAIS. Sans cette ligne le document est a 0 % de couverture, exactement
% comme le deck du 7 aout : c'est la faille qui avait laisse passer un intervalle faux. Le
% document n'ayant pas de \section, le harnais le lit comme un bloc unique, donc la citation se
% pose une fois, en fin de note, et porte sur l'ensemble.
{\scriptsize\noindent\textbf{Sources de cette note :} scripts~\texttt{15} et~\texttt{30}
(sensibilités et coût de l'ignorance directionnelle), \texttt{26} (protocole de Cooke),
\texttt{31} (effet d'un panel biaisé), \texttt{47} et~\texttt{63} (calibration figée),
\texttt{48} et~\texttt{71} (les trois étages), \texttt{69} (dépendance des états),
\texttt{70} (renormalisation par la taille), \texttt{72} (intervalle à $95\,\%$),
\texttt{73} (CTE agrégée), \texttt{74} (défaillances simultanées), \texttt{75} (séquences
ordonnées). Les grandeurs de la note sur $Y_{j,t}$ sont celles du script~\texttt{08b}.\par}

\vspace{0.2cm}
\begin{cle}
\textbf{Ce qui reste ouvert, et c'est court.} Deux lignes seulement appellent une décision :
le \textbf{niveau de l'intervalle reporté}, 90 ou $95\,\%$, qui ne déplace aucune calibration
mais change toutes les bornes publiées ; et le \textbf{renommage} \texttt{GBASE} contre
\texttt{G\_BASE} dans le code, qui touche 27 fichiers sur un pipeline gelé. Une troisième ligne
est non traitée et le restera faute de donnée, celle des dates d'arrivée. Tout le reste est soit
intégré, soit mesuré et écarté avec son motif.
\end{cle}

\end{document}
